\cotefile{157-3}
\scanpage{82}
\mspage{\unsure{31|34}}

Supposons p. ex. $f$ propre, considérons, \unsure{pour} $\deleted{\unsure{y}}s\in\mathrm{Ob}\,\unsure{S|\mathscr S}$
\diagram{\begin{tikzcd}[ampersand replacement=\&]
X_s \arrow[r,no head,"f_s"] \arrow[d,"i_s^X"'] \& Y_{\deleted{\unsure{y}}s} \arrow[d,"i_s^Y"]\\
X_{/\deleted{\unsure{y}}s} \arrow[r,no head,"f_{/s}"] \& Y_{/s},
\end{tikzcd}}
comme $X,Y$ propres, \deleted{\unsure{$X_s$} \illegible{1} \unsure{$Y_s$}} les flèches verticales sont dans $W$ (il suffirait de beaucoup moins que la propreté \ldots), donc $f_s\in W\Longrightarrow f_{/s}\in W$.

Ainsi, ``$f$ est dans $W$ loc. sur $S$'', ce qui implique $f\in W$ par $W(2)$.

\uline{Prop. 6} Soit
\diagram{\begin{tikzcd}[ampersand replacement=\&,column sep=small,row sep=small]
X \arrow[d,"f"'] \& X' \arrow[l,"{\unsure{g|q}}"'] \arrow[d,"f'"]\\
S \& S' \arrow[l,"p"]
\end{tikzcd}}
cartésien \unsure{dans|ds} Cat. Si $f$ $W$-asphérique, $p$ $W$-lisse, alors $f'$ est $W$-asph.

\deleted{\unsure{Dém.} Dire que $f$ est asphérique signifie ici que $f\colon X\to Y$ est dans $W$ \unsure{d'apr.}}

Dém. \deleted{En \unsure{utilisant} la propriété a) de prop. 1\deleted{, \unsure{et se}} pour \unsure{prendre} la notion de $W$-asphéricité, et se rappelant que les isom. locaux sont lisses et par suite si $Y''\to Y'$ est un \unsure{iso|iso.} local, $Y''$ est lisse sur $Y$, on est \unsure{ra}} Il faut prouver $X'\times_{S'}S'_{/s'}$ est $W$-asphérique ($s'\in\mathrm{Ob}\, S'$). Soit $y=p(y')$, on a
\diagram{\begin{tikzcd}[ampersand replacement=\&,column sep=small,row sep=small]
S \& S' \arrow[l,"p"'] \& {}\\
{} \& \deleted{\unsure{S}}S_{/s}\arrow[ul] \& S'_{/s'}\arrow[l,"p_0"]\arrow[ul],
\end{tikzcd}}
avec $p_0$ \inserted{$W$}-lisse \inserted{(\unsure{limites|lemme} \unsure{b.ses})} et à fibres $W$-asphériques, donc une $W$-équiv. universelle (prop. 5). Comme l'image inverse de $X$ sur $X_{/y}$ est $W$-asph. (hyp. sur $f$), il

% archive sheet number: 82
